% Данный файл распространяется под лицензией CC BY 4.0. Текст лицензии размещён на https://creativecommons.org/licenses/by/4.0/
% (c) Кафедра системного программирования, 2026

% Компилировать XeTeX

\documentclass[a4paper]{article}

\usepackage[a4paper, top=10mm, bottom=10mm, left=8mm, right=8mm]{geometry}

\usepackage{polyglossia}
\setdefaultlanguage[babelshorthands=true]{russian}
\setotherlanguage{english}

\usepackage{fontspec}
\setmainfont{FreeSerif}
\newfontfamily{\russianfonttt}[Scale=0.7]{FreeMono}

\usepackage[tiny, compact]{titlesec}

\usepackage{titling}
\setlength{\droptitle}{-1cm}
\pretitle{\begin{center}\begin{bfseries}\Large}
\posttitle{\par\end{bfseries}\end{center}}
\preauthor{\begin{center}\normalsize}
\postauthor{\par\end{center}\vspace{-1.8cm}}

\usepackage{hyperref}
\usepackage{bookmark}
\usepackage{enumitem}

\sloppy

\title{Обзор литературы к предрелизной подготовке системы MarLine}

\author{Овечкин Андрей}

\date{}

\begin{document}

\maketitle

\begin{flushright}
    Группа: \emph{25.Б72-мм}

    Кафедра: \emph{системного программирования}

    Научный руководитель: \emph{Гориховский Вячеслав Игоревич}

    Консультант: \emph{Дмитриевцев Алексей Сергеевич, младший инженер по разработке ПО, ООО ``КНС ГРУПП''}
\end{flushright}

\section{Введение}

Цель практики ~-- довести систему MarLine до состояния, при котором её можно
выпускать без ручной проверки: покрыть конвейер интеграционными тестами,
устранить найденные дефекты и привести репозиторий к публикуемому виду.
Прежде чем трогать код, нужно понять, на каких идеях система стоит, где она
повторяет уже опубликованные решения, а где выходит за их рамки. Без этого
не получится ни осмысленно чинить хешеры и индексы, ни объяснить, в чём
новизна работы. Ниже рассмотрены публикации, на которые MarLine опирается
архитектурно, а в конце сделаны выводы о
том, что из этого следует для предрелизной подготовки.

\section{Разбиение на чанки}

Дедупликация начинается с разбиения данных на блоки. Если резать по
фиксированному размеру, любая вставка в начало файла сдвигает все границы, и
неизменённая часть перестаёт совпадать с уже сохранённой копией. Поэтому в
современных системах используют разбиение, зависящее от содержимого: граница
ставится там, где свертка скользящего окна укладывается в заданное условие,
и положение границы определяется самими данными, а не порядковым номером
байта.

Полиномиальное хэширование на каждом байте дорого, поэтому FastCDC
\cite{fastcdc2016} заменил её Gear-хешем, который обновляется за один байт,
и добавил три приёма, избавляющих от лишней работы: упрощённую проверку
условия разреза, пропуск точек ближе минимального размера чанка и
нормализацию распределения размеров. Это дало примерно десятикратное
ускорение при том же качестве. Версия 2020-го года \cite{fastcdc2020} довела
проект до практики: приёмов стало пять, включая прокрутку окна на два байта,
ускорение выросло до 3--12$\times$, а на действующей системе Destor
(открытая платформа для экспериментов по дедупликации)
интеграция FastCDC ускорила сквозную запись в 1.2--3.0$\times$. Для нас
важно, что вторая статья проверяет алгоритм не в вакууме, а внутри настоящего
хранилища, ~-- она задаёт стандарт не только на сам чанкер, но и на то, как
его оценивать.

MarLine использует разбиение именно этой версии, совмещая его с отдельными
приёмами из статьи 2016 года. При предрелизной работе с
\texttt{fastcdc.rs} это означает, что таблицы Gear-констант и маски
выглядят взятыми из ниоткуда, но у каждой из них есть обоснование в одной из
двух статей, и при ревью кода удобно возвращаться к первоисточнику, а не
восстанавливать смысл по догадкам.

\section{Определение сходства}

Дедупликация распознаёт лишь идентичные блоки. Данные же резервных копий
часто похожи, но не совпадают: изменилась пара функций, файл пересобран
с другим набором флагов. Чтобы такие блоки сжимать, их нужно сначала найти,
а поиск похожих пар среди миллионов чанков ~-- самая дорогая часть работы.

Начиналась эта линия с системы Ароновича \cite{aronovich09}. Её метод сбора
признаков, известный как N-Transform, проходит хешем скользящего окна по
всему блоку и оставляет несколько наименьших значений в качестве признаков.
Авторы заметили, что распределение сверток неравномерно, и поправили отбор:
берутся четыре минимума, после чего поиск продолжается со сдвигом. Значение
этой работы для нас в том, что все последующие публикации измеряют себя
относительно N-Transform, ~-- это точка отсчёта, и её же мы видим в коде
как хешер сходства.

Дальше линия быстро расходится. ODESS \cite{odess2021} показал, что
просчитывать свертку на каждом окне не обязательно: устойчивый признак можно
получить из выборки, зависящей от содержимого, и применить преобразования
только к ней. Признаки стали генерироваться в 5.4$\times$ быстрее, чем у
Finesse, и в 26.9$\times$ быстрее, чем у N-Transform, при том же коэффициенте
сжатия. Finesse, промежуточный вариант, ускорял N-Transform ценой точности и
в итоге ниже фигурирует только как базовое сравнение.

Palantir \cite{palantir2024} предложил строить не один набор грубых
признаков, а иерархию: широкий слой находит как можно больше вариантов,
узкий отсеивает лишние. Сжатие бэкапа выросло на 7.3\% относительно
N-Transform и ODESS и на 26.5\% относительно Finesse при потере скорости не
более 7.7\%; дополнительные суперпризнаки авторы компенсировали тем, что
похожие данные в бэкапах приходят группами. Наконец, SpeedSketch
\cite{speedsketch2025} объединил поиск и кодирование в один конвейер: скетч
в стиле фильтра Блума не только ищет похожие блоки, но и сразу отбрасывает
те, которые всё равно не сожмутся. Генерация скетча в 1.59$\times$ быстрее,
чем у ODESS, кодирование в 1.92$\times$ быстрее Gdelta.

Для MarLine эта часть литературы важна ещё больше. Во-первых, система содержит
сразу две ветви сходства (ODESS-хешер в \texttt{marline\_sketcher} и
иерархические суперпризнаки в \texttt{marline\_palantir}), то есть внутри
одного проекта сосуществуют две линии из литературы, и их можно измерять на
общих данных. Во-вторых, именно здесь находится большинство найденных
дефектов: недетерминированные коэффициенты, из-за которых признаки двух
экземпляров хешера не совпадают, и ошибки в семантике индекса, из-за которых
повторные признаки искажают подсчет пересечений.

\section{Дельта-кодирование}

Найденную пару похожих блоков нужно записать компактно. Здесь состояние поля
удобно описать относительно Gdelta \cite{gdelta2020}: скорость ему даёт
Gear-хеш на мелких словах и массивная индексация, степень сжатия ~-- пакетная
обработка остатка после сопоставления слов. По замерам авторов, на шести
наборах данных кодирование и декодирование в 2--4$\times$ быстрее Xdelta и
Zdelta при росте коэффициента сжатия на 10--120\%.

Интересно, что в работах по определению сходства
\cite{odess2021,palantir2024,speedsketch2025} Gdelta выступает только как
база для сравнения и даже не цитируется. Это говорит о том, что задача
дельта-кодирования считается решённой, а исследовательский интерес лежит в
поиске похожих пар.

\section{Платформа}

Все перечисленные алгоритмы в MarLine живут не сами по себе, а внутри ChunkFS
\cite{chunkfs} ~-- in-memory файловой системы, где алгоритмы разбиения,
хеширования и хранения подменяются через трейты \texttt{Chunker} и
\texttt{Scrub}. Хдесь проходит граница между кодом:
всё, что относится к разбиению и хранению, система получает от платформы, а
собственные алгоритмы сходства и сжатия подключаются как скрабберы. Для
предрелизной подготовки это означает, что значительная часть конвейера ~--
внешняя зависимость, и её состояние нужно учитывать отдельно от состояния
собственного кода.

\section{Анализ}

Сведём рассмотренные методы определения сходства в таблицу по двум метрикам:
скорости генерации признаков и качеству сжатия. Приведены только значения,
указанные в источниках; там, где прямого числа нет, стоит прочерк.

\begin{center}
\begin{tabular}{l l l}
\hline
\textbf{Метод} & \textbf{Скорость признаков} & \textbf{Качество сжатия} \\
\hline
N-Transform & 1$\times$ (эталон) & 1$\times$ (эталон) \\
Finesse & --- & --- \\
ODESS & в 26.9$\times$ быстрее N-Transform & как у N-Transform \\
Palantir & медленнее N-Transform на $\le$7.7\% & на 7.3\% больше, чем у N-Transform\\
SpeedSketch & в 1.59$\times$ быстрее ODESS & кодирование в 1.92$\times$ быстрее Gdelta \\
\hline
\end{tabular}
\end{center}

Без прямых замеров остался только Finesse: в источниках есть лишь его
соотношения с другими методами (ODESS быстрее него в 5.4$\times$, Palantir
лучше него по сжатию на 26.5\%), поэтому в таблице он и помечен прочерком.

Из таблицы видно, что в литературе нет единого «лучшего» метода: ODESS
выигрывает скоростью, Palantir ~-- качеством, SpeedSketch ~-- скоростью
связки «скетч + дельта». Для MarLine это означает, что наличие двух ветвей
сходства -- ODESS-хешера и иерархических суперпризнаков Palantir -- не
избыточность, а два варианта решения проблемы. Их вклад возможно
оценить только экспериментально.

По разбиению на чанки выбора, по сути, нет: FastCDC 2020 ~-- практический
стандарт, и его использование в MarLine соответствует литературе. По
дельта-кодированию Gdelta ~-- устойчивый компромисс между скоростью и
качеством, и его выбор также обоснован.

\section{Выводы}

Из обзора следует три вывода.

Во-первых, архитектура MarLine повторяет проверенные решения: разбиение ~--
FastCDC 2020, сходство ~-- ODESS и Palantir, дельта ~-- Gdelta, платформа ~--
ChunkFS. Выбор обоснован литературой, менять его на этом этапе работы не
нужно.

Во-вторых, новизна работы не в отдельном алгоритме, а в интеграции и
проверке: в одной системе сосуществуют две ветви дедупликации на общей
платформе, и ни одна рассмотренная статья не даёт ни детерминированных
замеров, ни интеграционных тестов, падающих при ухудшении сжатия. Заполнение
этого пробела ~-- и есть предрелизная подготовка.

В-третьих, из анализа следует необходимость сравнить ветви ODESS и Palantir
экспериментально: литература даёт для каждой свои цифры, но на разных наборах
данных, и перенос итогов на MarLine без собственного замера некорректен.

\section{Следствия для постановки задачи}

Корректировка цели и задач практики по результатам обзора не требуется.
Сформулированная постановка ~-- интеграционное тестирование обеих ветвей,
устранение дефектов, проверка качества сжатия, приведение документации и CI
к публикуемому виду ~-- уже включает всё, что обзор показывает как пробел
в литературе. Таким образом, результаты обзора подтверждают постановку,
а не заставляют её изменять.

\begin{thebibliography}{9}

\bibitem{fastcdc2016}
W.\ Xia, Y.\ Zhou, H.\ Jiang, D.\ Feng, Y.\ Hua, Y.\ Hu, Y.\ Zhang, Q.\ Liu.
\textit{FastCDC: a Fast and Efficient Content-Defined Chunking Approach for Data
Deduplication}. Proceedings of USENIX Annual Technical Conference (ATC),
Denver, CO, 2016, pp.~101--114.

\bibitem{fastcdc2020}
W.\ Xia, X.\ Zou, H.\ Jiang, Y.\ Zhou, C.\ Liu, D.\ Feng, Y.\ Hua, Y.\ Hu,
Y.\ Zhang. \textit{The Design of Fast Content-Defined Chunking for Data
Deduplication Based Storage Systems}. IEEE Transactions on Parallel and
Distributed Systems, vol.~31 (9), 2020, pp.~2017--2031.

\bibitem{aronovich09}
L.\ Aronovich, R.\ Asher, E.\ Bachmat, H.\ Bitner, M.\ Hirsch, S.T.\ Klein.
\textit{The Design of a Similarity Based Deduplication System}. Proceedings of
SYSTOR~2009: The Israeli Experimental Systems Conference, Haifa, 2009,
pp.~1--14.

\bibitem{odess2021}
X.\ Zou, C.\ Deng, W.\ Xia, P.\ Shilane, H.\ Tan, H.\ Zhang, X.\ Wang.
\textit{Odess: Speeding up Resemblance Detection for Redundancy Elimination by
Fast Content-Defined Sampling}. 2021 IEEE 37th International Conference on
Data Engineering (ICDE), Chania, 2021, pp.~480--491.

\bibitem{palantir2024}
H.\ Huang, P.\ Wang, Q.\ Su, H.\ Xu, C.J.\ Xue, A.\ Brinkmann.
\textit{Palantir: Hierarchical Similarity Detection for Post-Deduplication
Delta Compression}. Proceedings of the 29th ACM International Conference on
Architectural Support for Programming Languages and Operating Systems
(ASPLOS), vol.~2, La Jolla, 2024, pp.~830--845.

\bibitem{speedsketch2025}
F.\ Yang, Y.\ Wang, C.\ Li, K.\ Zhou, H.\ Li. \textit{SpeedSketch: An
Ultra-Fast Sketch Generation and Delta Encoding Framework for Delta
Compression}. Proceedings of the 54th International Conference on Parallel
Processing (ICPP), San Diego, 2025, pp.~168--177.

\bibitem{gdelta2020}
H.\ Tan, Z.\ Zhang, X.\ Zou, Q.\ Liao, W.\ Xia. \textit{Exploring the
Potential of Fast Delta Encoding: Marching to a Higher Compression Ratio}.
2020 IEEE International Conference on Cluster Computing (CLUSTER), Kobe, 2020,
pp.~198--208.

\bibitem{chunkfs}
O.\ Piletskii. \textit{ChunkFS} ~-- файловая система для сравнения алгоритмов
дедупликации. \url{https://github.com/Piletskii-Oleg/chunkfs}

\end{thebibliography}

\end{document}
